% =====================================================================
\section{From Object to Pixel: View, Projection, Clipping, Rasterisation, Depth}\label{sec:geomchain}
\pipestrip{3}
A vertex has left its object ($\vect p_\text{world}=W\,\vect p_\text{local}$, Chapters~\ref{sec:builder}--\ref{sec:transform}). This chapter follows it until it is a pixel with a depth:
\begin{equation}
\vect p_\text{clip}=P\,V\,W\,\vect p_\text{local},\qquad \vect p_\text{ndc}=\frac{\vect p_\text{clip}}{w_\text{clip}},\qquad
\begin{pmatrix}x_s\\y_s\\z_s\end{pmatrix}=\begin{pmatrix}(x_\text{ndc}+1)\,W_s/2\\(y_\text{ndc}+1)\,H_s/2\\(z_\text{ndc}+1)/2\end{pmatrix}.
\end{equation}
\emph{Reading: world, then camera frame, then squashed by the lens into a cube, divided by depth, stretched onto the screen.}

\subsection{View matrix: lookAt}
The camera sits at $\vect e$ looking at $\vect c$ with up $\vect u_0$. Its frame is $\vect f=\mathrm{normalize}(\vect c-\vect e)$, $\vect s=\mathrm{normalize}(\vect f\times\vect u_0)$, $\vect u=\vect s\times\vect f$ (right, true up, forward), and the view matrix expresses world points in it:
\begin{equation}
V=\begin{pmatrix}s_x&s_y&s_z&-\vect s\!\cdot\!\vect e\\u_x&u_y&u_z&-\vect u\!\cdot\!\vect e\\-f_x&-f_y&-f_z&\vect f\!\cdot\!\vect e\\0&0&0&1\end{pmatrix}.
\end{equation}
\emph{Reading: translate the eye to the origin, then rotate so that the camera looks down $-z$.} (\code{lookAt}, Listing~\ref{lst:LookAt}.)
\lstLookAt

\subsection{Projection matrix: perspective}
With vertical field of view $\theta$, aspect ratio $a$ and near and far planes $n,f$, and $t=1/\tan(\theta/2)$:
\begin{equation}
P=\begin{pmatrix}t/a&0&0&0\\0&t&0&0\\0&0&-\dfrac{f+n}{f-n}&-\dfrac{2fn}{f-n}\\0&0&-1&0\end{pmatrix}.
\end{equation}
\emph{Reading: $x$ and $y$ are scaled by the focal length; the last row copies $-z_\text{view}$ into $w$, so the later division by $w$ makes distant things smaller.} Depth is non-linear: for $d=-z_\text{view}$,
$z_\text{ndc}=\dfrac{f+n}{f-n}-\dfrac{2fn}{(f-n)\,d}$, so $d=n\mapsto-1$ and $d=f\mapsto+1$; most of the range is spent near the camera, so the depth view linearises it, $z_\text{lin}=2nf/(f+n-z_\text{ndc}(f-n))$. The orthographic matrix of the moon's shadow cascades has no division: $2/(r-l)$, $2/(t-b)$, $-2/(f-n)$ on the diagonal plus a translation column.
\lstPerspective
\paragraph{Worked example.} Take $\theta=60^\circ$ ($t=1.7321$), $a=16/9$, $n=0.1$, $f=200$ and a point at view-space $(1,\,0.5,\,-10)$. Then $x_c=1.7321/1.7778\cdot1=0.9743$, $y_c=1.7321\cdot0.5=0.8660$,
$z_c=\frac{200.1}{199.9}\cdot10-\frac{40}{199.9}=9.810$, $w_c=10$. Dividing by $w_c$: $\vect p_\text{ndc}=(0.0974,\,0.0866,\,0.9810)$, inside $[-1,1]^3$, so it is visible. On a $1920\times1080$ viewport the pixel is $x_s=(0.0974+1)\cdot960=1053.5$,
$y_s=(0.0866+1)\cdot540=586.8$, window depth $z_s=0.9905$.

\subsection{Clipping, NDC and the viewport}
In clip space a point is inside if $-w\le x,y,z\le w$; the hardware clips triangles against these six planes \emph{before} the division (avoiding $w\le0$). Afterwards the visible volume is the cube $[-1,1]^3$ (NDC), and the viewport transform maps it to pixels. The program also culls whole parts: a part's bounding sphere
$r=\tfrac12\sqrt{s_x^2+s_y^2+s_z^2}$ is tested against the six planes of $PV$ (Gribb--Hartmann: plane $=\text{row}_4\pm\text{row}_{1,2,3}$).

\subsection{Back-face culling and rasterisation}
After projection the signed area $\tfrac12[(x_1-x_0)(y_2-y_0)-(x_2-x_0)(y_1-y_0)]$ is positive for front faces; negative ones are discarded (\code{GL\_CCW}, \code{GL\_BACK}). The \emph{rasteriser} then finds the pixels whose centres lie inside the triangle with edge functions
$E_{01}(\vect p)=(x_1-x_0)(y-y_0)-(y_1-y_0)(x-x_0)$ (inside if the three have the same sign), and the barycentric coordinates are the normalised edge functions: $\lambda_2=E_{01}(\vect p)/2A$, and so on ($A$ the triangle area). These $\lambda$ are what Chapter~\ref{sec:interp} interpolates with.
\begin{figure}[h]
\centering
\begin{subfigure}{0.49\linewidth}\picdia{raster_edges}{\linewidth}\caption{edge functions and barycentric weights}\end{subfigure}
\begin{subfigure}{0.49\linewidth}\picdia{frustum_ndc}{\linewidth}\caption{view frustum $\to$ NDC cube}\end{subfigure}
\caption{Rasterisation and the view volume.}
\end{figure}

\subsection{Depth test: the Z-buffer}
Each fragment's depth $z_s$ is compared with the stored value (\code{GL\_LESS}): if smaller, colour and depth are written, otherwise the fragment is discarded. Key \texttt{Z} switches the test off, showing the painter's-order result (e.g.\ the rover through the tree that hides it). The depth buffer is reused by the shadow maps, SSAO, fog, SSR and the vision's ground-ray hits.
\begin{figure}[h]
\centering
\piclab{buffer_03_depth}{0.7\linewidth}
\caption{The depth buffer of the current scene, linearised for display (in-program render, near = dark).}\label{fig:diag}
\end{figure}

\subsection{The workbook view of the same matrices}
The workbook rebuilds $T$, $R$, $S$ and $M=TRS$ of every part with sum-of-products cell formulas ($m_{ij}=\sum_k a_{ik}b_{kj}$) and checks them against the exported matrices: the largest difference is $8.7\times10^{-8}$ on the UAV gimbal chain and $2.4\times10^{-5}$ over the whole workbook (single-precision rounding).
\begin{figure}[h]
\centering
\picsheet{sheet_mmult}{\linewidth}
\caption{\sheetcaption{sheet_mmult}}
\end{figure}
